%%%%%%
%
% $Autor: Wings $
% $Datum: 2019-03-05 08:03:15Z $
% $Pfad:  $
% $Version: 4250 $
% !TeX spellcheck = en_GB
% !TeX encoding = utf8
% !TeX root = ArduinoIDE2Win 
% !TeX TXS-program:bibliography = txs:///biber
%
%%%%%%


\subsection{Installation on Windows}

The installation process for Arduino IDE 2 on a Windows computer will be described step by step in the upcoming section. Following the installation instructions, the subsequent chapter will provide an overview of the IDE's features, including the Board Manager, Library Manager, Sketchbook, and more.

To install the Arduino IDE 2 on a Windows computer, simply run the file downloaded from the software page Arduino Software \HREF{https://www.arduino.cc/en/software}{Arduino-Software}. Follow the installation steps, as shown in the Figure ~\ref{downloading-and-installing-img02}, to ensure correct installation.


\begin{center}\centering
	\includegraphics[width=12cm]{Arduino/ArdiunoIDE2/WindowsDownload.png}
	\captionof{figure}{Arduino IDE Create Agent Installation}\label{fig:ArduinoIDECreatAgentInstallation}		
\end{center}




After the download is complete, open the \menu {file setup} and proceed with the installation. Select all components as shown in Figure ~\ref{fig:ArduinoSetupInstallationFolderoptions} in the dialog box, and then click \menu {Next}
.


\begin{center}\centering
	\includegraphics[width=8cm]{Arduino/ArdiunoIDE2/ArduinoSetupInstallationOptions}
	\captionof{figure}{Arduino Setup Installation options}\label{fig:ArduinoSetupInstallationFolderoptions}		
\end{center}

Select the destination folder as seen in Figure ~\ref{fig:ArduinoSetupInstallationFlder} and click \menu {Install}.

\begin{center}\centering
	\includegraphics[width=8cm]{Arduino/ArdiunoIDE2/ArduinoSetupInstallationFolder}
	\captionof{figure}{Arduino Setup Installation Folder}\label{fig:ArduinoSetupInstallationFlder}		
\end{center}

Once the installation is complete, open the Arduino IDE. A default sketch will appear on the screen, as shown in Figure ~\ref{fig:ArduinoIDESketch}.

\begin{center}\centering
	\includegraphics[width=8cm]{Arduino/ArdiunoIDE2/ArduinoIDESketch}
	\captionof{figure}{Arduino IDE Sketch}
	\label{fig:ArduinoIDESketch}		
\end{center}


It can be seen from the above figure~\ref{fig:ArduinoIDESketch} that the basic arduino sketch has two parts. The first part is the function \PYTHON{void setup()} which returns void and we do the intiliaztion such as the output LED color, specifying the core etc. The second part is the function \PYTHON{void loop()} where we define functions which are to be performed through out the loop. These codes are placed between paranthesis \PYTHON{$\{ \}$} and each function has a return type, here it has void return type.


Below in Figure~\ref{DownloadingAndInstallingImg02} is a summary of the installation steps for the Windows version, outlining the key actions required to properly install the softwar

\begin{center}\centering
	
	\includegraphics[width=\linewidth]{Arduino/ArdiunoIDE2/DownloadingAndInstallingImg02.png}
	\captionof{figure}{Steps for Installation}\label{DownloadingAndInstallingImg02}
	
\end{center}

